\section{Checks with GANDALF}\label{sec:numerics}

We checked the results of sections~\ref{sec:slices} to~\ref{sec:truncated} with the GANDALF code \citep[version X;][]{Kanekar2025}, % GANDALF release version: to be filled by the author
which integrates (\ref{eq:rmhd}) with the truncated Hermite hierarchy of one branch (section~\ref{sec:model}). All computations are in double precision on a CPU, the GANDALF runs use its fourth-order Lawson--Runge--Kutta time step (Lawson-RK4), and we computed every diagnostic from the saved moments and potentials, not with the code's built-in routines. The box is periodic, with side $L=1$ and $k_1=2\pi/L$. As section~\ref{sec:truncated} requires, the fastest slice sets the time step: $\Delta t$ is chosen with a safety factor of one half on the smaller of the Alfv\'enic limit and $2\sqrt2/\rho$, the stability limit of RK4 on the imaginary axis, with $\rho=k_{\perp,\max}\max_j|\nabla_\perp(\Phi+su_j\Psi)|$. The configuration files and scripts accompany the paper.

\subsection{Spectral density and phase velocities}\label{num:spectrum}

\begin{figure}[tbp]
  \centering
  \includegraphics[width=\textwidth]{spectrum}
  \caption{(a) The spectral density $\varrho(u)$ of (\ref{eq:varrho}) for $\Lambda=\sqrt5$, $-\sqrt5$ and $\infty$ (lines), and the Gauss weights $\lambda_j$ of GANDALF's ladder matrix at $M=16$ with the zero closure, divided by the node spacing $\Delta u_j=(u_{j+1}-u_{j-1})/2$, at the nodes $u_j$ (markers in the colour of each line). (b) The largest phase velocity $u_{\max}$ of GANDALF's ladder matrix at $\Lambda=\sqrt5$, minus $(2M+3)^{1/2}$, for $M=4$ to $1024$ with the zero closure (circles) and the copy closure (squares), and the three-term (solid) and two-term (dashed) asymptotic forms (\ref{eq:umax}).}
  \label{fig:spectrum}
\end{figure}

In figure~\ref{fig:spectrum}(a), since $\eps(0)=p_0$, $\varrho(0)=F_M(0)/p_0^2$ is $1.8463$, $0.2694$ and $0.5642$ for $\Lambda=\sqrt5$, $-\sqrt5$ and $\infty$, and $\int\varrho\,\dd u=1/p_0$ to round-off. By (\ref{eq:gauss}), $\lambda_j/\Delta u_j$ approximates $\varrho(u_j)$; for GANDALF's ladder matrix at $M=16$ the two differ on $|u_j|<3$ by at most 8.2\%, 1.8\% and 0.47\%, respectively. In figure~\ref{fig:spectrum}(b) the three-term forms (\ref{eq:umax}) follow the largest eigenvalue of GANDALF's ladder matrix down to $M=4$ for both closures; at $M=64$ and 128 it has the values quoted in section~\ref{sec:truncated}, and it does not depend on $\Lambda$ to round-off.

\subsection{Slices in a turbulent run}\label{num:snapshots}

\begin{figure}[tbp]
  \centering
  \includegraphics[width=\textwidth]{snapshots}
  \caption{A density blob in decaying Alfv\'enic turbulence, on the plane $z=0$. GANDALF run on a $64^3$ grid with $M=16$ (zero closure), $\Lambda=\sqrt5$ and $\beta_i=1$; balanced Gaussian random Elsasser fields with $1\le k_\perp/k_1\le3$, $|k_z|/k_1\le2$, energy-weighted mean wavenumbers $\bar k_\perp=2.45\,k_1$ and $\bar k_\parallel=1.17\,k_1$, and $z_{\rm rms}=0.478$, the rms of $|\nabla_\perp z^\pm|$, so that $\chi=\bar k_\perp z_{\rm rms}/\bar k_\parallel=1.00$ (0.73 with the rms of $|\uE|$); hyper-resistivity damps $z^\pm$ and every $g_m$ at the rate $2(k_\perp/k_{\perp,\max})^4$; no forcing or hypercollisions; $g(0)=f\,e_0$, with $f$ a Gaussian of standard deviation $L/8$ in $x$ and $y$ and $L/4$ in $z$; 148 steps to $t=1.5\,t_{\rm nl}$, where $t_{\rm nl}=1/(\bar k_\perp z_{\rm rms})=0.1357$ is the outer eddy time. (a, b) The density moment $g_0$ at $t=0$ and $1.5\,t_{\rm nl}$. (c--e) The slices $a_j/(p_0R_{0j})$ at $t=1.5\,t_{\rm nl}$ at the nodes $su_j=-1.0218$, $u_j=0$ and $su_j=+1.0218$; all three equal $f$ at $t=0$. Each group of panels has one symmetric colour scale.}
  \label{fig:snapshots}
\end{figure}

The run of figure~\ref{fig:snapshots} releases a density blob, $g(0)=f\,e_0$, into decaying Alfv\'enic turbulence with $\chi=1.00$. Its hyper-resistivity acts alike on $z^\pm$ and every $g_m$, so by (\ref{eq:sliceforcing}) it damps each slice separately. Since $a_j=\sum_mp_mR_{mj}g_m$, every slice starts as $p_0R_{0j}f$, and by (\ref{eq:truncslice}) each is then carried by its own stream function $\Phi+su_j\Psi$, which at $su_j=-1.0218$, $0$ and $+1.0218$ is mainly $z^-$, $\Phi$ alone and mainly $z^+$ (proposition~\ref{prop:elsasser}). These three slices start from the same pattern and end visibly different, and $g_0=\sum_jR_{0j}a_j$ mixes them.

At $t=0$, $0.5\,t_{\rm nl}$ and $1.5\,t_{\rm nl}$, GANDALF's $\partial_tg$ projected on the slices agrees with $-\{\Phi+su_j\Psi,a_j\}-su_j\partial_za_j$, computed with brackets that do not use GANDALF, to $4.8\times10^{-16}$, $5.2\times10^{-16}$ and $5.2\times10^{-16}$ of the largest slice rate (per slice, up to $4.5\times10^{-12}$ at the extreme nodes, which hold about $10^{-13}$ of the energy of the central slice), while the bracket term is 0.56 to 2.2 times the streaming term in norm. Advecting every slice by $\Phi$ alone instead gives relative errors of 0.30 to 1.46, except at $u_j=0$, where the error is $5.6\times10^{-16}$. Over $1.5\,t_{\rm nl}$ the slice energies $\int a_j^2\,\dd^3r$ change by at most $4.3\times10^{-4}$ of their initial values, and each change equals the hyper-resistive loss of that slice to $1.2\times10^{-6}$ of its initial energy (the total to $6\times10^{-12}$), so no free energy passes between slices. At the end, $k_\perp/k_1\ge21$ holds $8.9\times10^{-7}$ of the free energy and $1.2\times10^{-5}$ of the Alfv\'enic energy.

\subsection{The Hermite spectrum of decorrelated slices}\label{num:dictionary}

\begin{figure}[tbp]
  \centering
  \includegraphics[width=\textwidth]{dictionary}
  \caption{(a) The Hermite spectrum $W(m)$ in units of its large-$m$ form $\int e\,\dd u/(\pi\sqrt{2m})$. Lines: $\delta$-correlated slices, with the closed form (\ref{eq:prefactor}) for $e\propto F_M$ and with (\ref{eq:dictionary}), evaluated by Gauss quadrature, for $e\propto\varrho$ at $\Lambda=\sqrt5$ and $-\sqrt5$. Markers: the steady state of one Fourier mode whose $g_0$ is driven by white noise, for $\Lambda=\infty$, $\sqrt5$ and $-\sqrt5$, from GANDALF's ladder matrix at $M=600$ with the outgoing closure $g_{M+1}=-\ii\,\mathrm{sgn}(k_z)\,g_M$ and the Lyapunov equation for the steady covariance. (b) The local slope $\dd\ln W/\dd\ln m$ from (\ref{eq:dictionary}) for correlated slices with $e=F_M$ and $\hat c(\kappa)=(1+\kappa^2\ell^2)^{-p/2}$, for $p=0.5$ to $2.5$ with $\ell=1$ (solid) and $\ell=0.15$ (dashed); the dotted lines are $-(p+1)/2$.}
  \label{fig:dictionary}
\end{figure}

By section~\ref{sec:decorrelation}, the forced linear steady state (markers in figure~\ref{fig:dictionary}a) is the spectrum of $\delta$-correlated slices with $e\propto\varrho$ (lines). At $\Lambda=\infty$ it differs from the closed form (\ref{eq:prefactor}) by at most $3.1\times10^{-4}$, and at $\Lambda=\sqrt5$ and $-\sqrt5$ from the spectrum with $e\propto\varrho$ by at most $1.5\times10^{-4}$ and $4.2\times10^{-4}$, the size of the closure error. At $m=1$ it is 4.7\% below $(1-1/\Lambda)$ times its value at $\Lambda=\infty$ for $\Lambda=\sqrt5$ and 5.2\% above it for $\Lambda=-\sqrt5$, so this proportionality holds only at large $m$. In figure~\ref{fig:dictionary}(b) the local slope is within $3\times10^{-5}$ of $-(p+1)/2$ at $m=10^4$ for $\ell=1$; for $\ell=0.15$ it reaches these values at larger $m$, where $\sqrt{2m}\,\ell\gg1$.

\subsection{Frozen field lines}\label{num:frozen}

\begin{figure}[t]
  \centering
  \includegraphics[width=\textwidth]{frozen}
  \caption{Test of the frozen-line theorem. Ideal GANDALF runs (no dissipation, forcing or collisions) from one initial state at $N^3$ grid points, with $M=8$ (zero closure), $\Lambda=\sqrt5$, $\beta_i=1$, $\mathrm{rms}\,|\nabla_\perp\Phi|=\mathrm{rms}\,|\nabla_\perp\Psi|=1$, random $g_0$ and $g_1$, $\chi=2.21$ (1.56 with the rms of $|\uE|$), the term $\{\Phi,g_0\}$ 2.91 times the streaming term $s\partial_z(Ag)_0$ in rms, and $\Delta t=0.0005$. (a) The Hermite spectrum $W(m)$ at $t=0.125$ for $N=64$: initial state (crosses), GANDALF (circles), frozen-line prediction (solid line) and $k_z$-only prediction (dashed line). (b) The relative difference between GANDALF and the frozen-line prediction at $t=0.125$ against $N$, for the single-point statistics $W(m)$, $C_{jl}$, $\Pi_{m+1/2}$ with $\gradpar$ and $\langle g_m^4\rangle$ (filled markers), and for $W(k_\perp)$, $W(|k_z|)$ and $\Pi_{m+1/2}$ with $\partial_z$ (open markers). Table~\ref{tab:frozen} gives the values.}
  \label{fig:frozen}
\end{figure}

\begin{table}[t]
  \centering
  \caption{Relative differences between GANDALF ($X$) and the frozen-line prediction ($X_{\rm p}$) for the runs of figure~\ref{fig:frozen} at $t=0.062$ and $0.125$, measured by $\sum|X-X_{\rm p}|/\sum|X_{\rm p}|$ for $W(m)$, $\langle g_m^4\rangle$ ($m\le2$), $W(k_\perp)$ and $W(|k_z|)$ and by $\max|X-X_{\rm p}|/\max|X_{\rm p}|$ for $C_{jl}$ and the fluxes; the last two rows of each block use the $k_z$-only prediction. At $t=0.062$ the values of $W(m)$, $C_{jl}$ and $\langle g_m^4\rangle$ at $N=64$ are limited by the time step.}
  \label{tab:frozen}
  \small
  % Generated by analysis/fig_frozen.py from runs P1_frozen-N16_20261014_152357, P1_frozen-N20_20261014_152405, P1_frozen-N32_20261014_152417, P1_frozen-N40_20261014_152443, P1_frozen-N64_20261014_152540.
% Relative differences between GANDALF (gandalf_step, Lawson-RK4) and the frozen-line prediction (last two rows: the k_z-only prediction).
\begin{tabular}{lccccc}
\toprule
statistic & $N=16$ & $N=20$ & $N=32$ & $N=40$ & $N=64$ \\
\midrule
\multicolumn{6}{l}{$t=0.062$} \\
\midrule
$W(m)$ & $4.5\times10^{-4}$ & $2.1\times10^{-4}$ & $1.3\times10^{-5}$ & $2.4\times10^{-6}$ & $2.3\times10^{-7}$ \\
$C_{jl}=\langle a_ja_l\rangle$ & $2.5\times10^{-4}$ & $1.8\times10^{-4}$ & $1.1\times10^{-5}$ & $1.9\times10^{-6}$ & $1.4\times10^{-7}$ \\
$\Pi_{m+1/2}$ with $\nabla_\parallel$ & $4.1\times10^{-3}$ & $2.1\times10^{-3}$ & $1.6\times10^{-4}$ & $3.3\times10^{-5}$ & $1.1\times10^{-6}$ \\
$\langle g_m^4\rangle$, $m\le2$ & $4.9\times10^{-3}$ & $2.1\times10^{-3}$ & $5.2\times10^{-5}$ & $4.2\times10^{-6}$ & $5.2\times10^{-7}$ \\
\addlinespace
$W(k_\perp)$ & $0.21$ & $0.21$ & $0.21$ & $0.21$ & $0.21$ \\
$W(|k_z|)$ & $0.15$ & $0.15$ & $0.15$ & $0.15$ & $0.15$ \\
$\Pi_{m+1/2}$ with $\partial_z$ & $0.069$ & $0.070$ & $0.069$ & $0.069$ & $0.069$ \\
\addlinespace
$W(m)$, $k_z$-only prediction & $0.077$ & $0.078$ & $0.078$ & $0.078$ & $0.078$ \\
$\Pi_{m+1/2}$, $k_z$-only prediction & $0.44$ & $0.44$ & $0.45$ & $0.45$ & $0.45$ \\
\midrule
\multicolumn{6}{l}{$t=0.125$} \\
\midrule
$W(m)$ & $3.2\times10^{-3}$ & $1.3\times10^{-3}$ & $2.8\times10^{-4}$ & $1.5\times10^{-4}$ & $1.5\times10^{-5}$ \\
$C_{jl}=\langle a_ja_l\rangle$ & $2.8\times10^{-3}$ & $1.0\times10^{-3}$ & $3.8\times10^{-4}$ & $1.1\times10^{-4}$ & $2.2\times10^{-5}$ \\
$\Pi_{m+1/2}$ with $\nabla_\parallel$ & $0.017$ & $5.3\times10^{-3}$ & $1.5\times10^{-3}$ & $7.7\times10^{-4}$ & $1.4\times10^{-4}$ \\
$\langle g_m^4\rangle$, $m\le2$ & $0.020$ & $0.012$ & $1.6\times10^{-3}$ & $2.0\times10^{-4}$ & $5.5\times10^{-5}$ \\
\addlinespace
$W(k_\perp)$ & $0.38$ & $0.38$ & $0.38$ & $0.38$ & $0.38$ \\
$W(|k_z|)$ & $0.26$ & $0.26$ & $0.26$ & $0.27$ & $0.27$ \\
$\Pi_{m+1/2}$ with $\partial_z$ & $0.13$ & $0.13$ & $0.14$ & $0.14$ & $0.14$ \\
\addlinespace
$W(m)$, $k_z$-only prediction & $0.15$ & $0.15$ & $0.15$ & $0.15$ & $0.15$ \\
$\Pi_{m+1/2}$, $k_z$-only prediction & $0.37$ & $0.37$ & $0.37$ & $0.37$ & $0.37$ \\
\bottomrule
\end{tabular}

\end{table}

The runs of figure~\ref{fig:frozen} integrate one initial state, drawn on an $8^3$ grid and embedded, at $N=16$, 20, 32, 40 and 64, to $t=0.25$, about three eddy times. With $\Delta t=0.0005$, $\rho\,\Delta t\le1.51$; at $t=0.25$ the free energy and the Alfv\'enic energy drift by at most $2.6\times10^{-8}$ and $4.3\times10^{-5}$, and halving $\Delta t$ divides the former by 32. We compare each run with the frozen-line prediction $h(t)=e^{-tL_0}g(0)$, computed with GANDALF's right-hand side for $g$ with $\Phi=0$ and $\Psi=\Psi(0)$ held fixed, and with the $k_z$-only prediction ($\Phi=\Psi=0$). By corollary~\ref{cor:singlepoint}, $W(m)$, $C_{jl}=\langle a_ja_l\rangle$, the flux (\ref{eq:fluxCD}) with $\gradpar(t)$ in the run and $\gradpar(0)$ in the prediction, and $\langle g_m^4\rangle$ match the frozen-line prediction up to the Fourier truncation error, which vanishes as $N$ grows at fixed $t$; $W(k_\perp)$, $W(|k_z|)$ and the flux with $\partial_z$ are not covered.

At $t=0.125$ the relative difference falls from $3.2\times10^{-3}$ at $N=16$ to $1.5\times10^{-5}$ at $N=64$ for $W(m)$, from $2.8\times10^{-3}$ to $2.2\times10^{-5}$ for $C_{jl}$, from $1.7\times10^{-2}$ to $1.4\times10^{-4}$ for the flux and from $2.0\times10^{-2}$ to $5.5\times10^{-5}$ for $\langle g_m^4\rangle$. It stays at 0.38 for $W(k_\perp)$, 0.26--0.27 for $W(|k_z|)$ and 0.13--0.14 for the flux with $\partial_z$ at every $N$. The $k_z$-only prediction misses $W(m)$ by 15\% and the field-aligned flux by 37\%. At $N=64$ halving $\Delta t$ changes these values by less than 2\%. At every $N$ the difference grows from $t=0.062$ to $0.125$ (table~\ref{tab:frozen}), as the Fourier truncation error is not uniform in $t$, which restricts such tests to times at which the Alfv\'enic fields stay resolved.

The scripts in \texttt{derivations/dynamics/} test two further statements of section~\ref{sec:lagrangian} in the same way. In ideal turbulence forced through $\Phi$, the response of $W(m)$ to each white injection into $g$ converges with $N$ to phase mixing along the field lines frozen at its injection time, and the average over the signs of the injections is the sum of these responses, as corollary~\ref{cor:forcing} states, although single realisations contain cross terms. With one term added, the departure from the frozen-line prediction still falls with $N$ for viscosity or forcing of $\Phi$, but not for resistivity or forcing of $\Psi$ or for damping of $g$ in $k_\perp$ or $k_z$.

\subsection{Further checks and the scripts}\label{num:rhs}

The other checks cited in sections~\ref{sec:model} to~\ref{sec:truncated} use GANDALF's right-hand side. For random states with $M=2$ to 65, $\Lambda=\pm\sqrt5$ and 3, $\beta_i$ from 0.37 to 1.7 and both closures, its $\partial_tg$ satisfies (\ref{eq:truncslice}) to $2\times10^{-13}$, with brackets several times the streaming term, and fails at order one when the sign of $\{\Psi,\cdot\}$ is reversed. Slices at $su_j=-1$ obey proposition~\ref{prop:elsasser} to $5.5\times10^{-16}$ and $1.6\times10^{-15}$, and translate rigidly when $z^-=0$, to $1.4\times10^{-14}$. The two-slice law (\ref{eq:twoslice}) holds to $5\times10^{-15}$ and the two forms of (\ref{eq:fluxCD}) agree to $5\times10^{-14}$. The forms $Q^{(M)}_\varphi$ are conserved to relative rates of $3.2\times10^{-17}$, whereas $\int a_j^4\,\dd^3r$ changes at relative rates of $4.9\times10^{-4}$ and more, so only the quadratic invariants of section~\ref{sec:invariants} survive Fourier truncation.

The self-checking scripts accompanying this paper are in \texttt{derivations/}: \texttt{continuum/} (sections~\ref{sec:model} to~\ref{sec:invariants}, appendices~\ref{app:branches} and~\ref{app:spectral}), \texttt{truncated/} (sections~\ref{sec:decorrelation} and~\ref{sec:truncated}, appendix~\ref{app:truncated}), \texttt{dynamics/} (sections~\ref{sec:slices}, \ref{sec:decorrelation} and~\ref{sec:lagrangian}), and \texttt{review/}, with independent re-derivations that share no code with the other folders. Each script prints PASS or FAIL for every check; a folder runs from the repository root with
\begin{center}
  \texttt{bash paper/phase-velocity-slices/derivations/}\textit{folder}\texttt{/run\_all.sh}
\end{center}
The figures are produced by the scripts in \texttt{scripts/} and \texttt{analysis/} from the configuration files in \texttt{configs/}.
